%!TEX root = ../main.tex
\chapter{Klasyczna infrastruktura klucza publicznego}
\label{ch:pki-klasyczna}

\section{Kryptografia asymetryczna i podpisy cyfrowe}
\label{sec:kryptografia-asymetryczna}

\section{Standard X.509 i certyfikaty}
\label{sec:x509}

\section{Hierarchia zaufania i rola urzędów certyfikacji}
\label{sec:hierarchia-zaufania}

\section{Mechanizmy unieważniania (CRL, OCSP)}
\label{sec:crl-ocsp}

\section{Słabości modelu scentralizowanego}
\label{sec:slabosci-pki}

\subsection{Pojedynczy punkt zaufania}
\label{subsec:spof}

\subsection{Studia przypadków: DigiNotar i Symantec}
\label{subsec:case-studies}

\section{Certificate Transparency jako częściowa odpowiedź}
\label{sec:ct}